\documentclass[11pt]{article}
\usepackage[margin=1in]{geometry}
\usepackage[T1]{fontenc}
\usepackage{mathptmx}\usepackage[scaled=.92]{helvet}\usepackage{courier}\usepackage{microtype}\renewcommand{\familydefault}{\sfdefault}
\usepackage{booktabs,tabularx,array,xcolor,graphicx,enumitem,fancyhdr,titlesec,listings,amsmath}
\usepackage[hidelinks]{hyperref}
\definecolor{accent}{HTML}{1F4E79}
\definecolor{soft}{HTML}{EEF3F8}
\definecolor{good}{HTML}{2E7D32}
\titleformat{\section}{\large\bfseries\color{accent}}{\thesection}{0.6em}{}
\titleformat{\subsection}{\normalsize\bfseries\color{accent}}{\thesubsection}{0.6em}{}
\setlength{\parindent}{0pt}\setlength{\parskip}{0.55em}
\renewcommand{\arraystretch}{1.25}
\pagestyle{fancy}\fancyhf{}
\fancyfoot[C]{\small\thepage}
\renewcommand{\headrulewidth}{0pt}
\lstset{basicstyle=\ttfamily\small,backgroundcolor=\color{soft},frame=none,breaklines=true,columns=fullflexible,keepspaces=true,xleftmargin=0.5em,aboveskip=0.6em,belowskip=0.4em}
\newcommand{\code}[1]{\texttt{\small\detokenize{#1}}}
\usepackage[most]{tcolorbox}
\newtcolorbox{callout}[1]{colback=soft,colframe=accent,boxrule=0.5pt,arc=2pt,left=6pt,right=6pt,top=4pt,bottom=4pt,title={#1},fonttitle=\bfseries}
\usepackage{colortbl}
\newcommand{\titleblock}[2]{\begin{center}{\LARGE\bfseries\color{accent}#1}\\[0.4em]{\large #2}\\[0.3em]{\small Kevin Williams \quad|\quad October 23, 2026}\end{center}\vspace{0.3em}}
\begin{document}
\titleblock{Verification Process}{A sample procedure showing the volatility surface tool works as intended}

\section{Approach: compare against answers known in advance}
A plot that looks plausible is not proof. This procedure checks the tool against answers that are known \emph{before} it runs, using four kinds of evidence of rising strength: known formulas, round trips, a planted truth, and independent programs that share no code with MATLAB. It takes about two minutes and needs no market data for the first six checks.

\begin{callout}{Run it}
In MATLAB: \code{verify_demo} \quad offline checks only\\[0.3em]
In MATLAB: \code{verify_demo("snapshots/<file>.csv")} \quad adds the real-data check\\[0.3em]
In a terminal: \code{python3 verify_independent.py snapshots/<file>.csv} \quad independent check
\end{callout}

\section{The checks}
\begin{tabularx}{\linewidth}{@{}c>{\raggedright\arraybackslash}p{0.27\linewidth}>{\raggedright\arraybackslash}X>{\raggedright\arraybackslash}p{0.2\linewidth}@{}}
\toprule
\textbf{\#} & \textbf{Check} & \textbf{Why it proves something} & \textbf{Pass if}\\\midrule
1 & Textbook prices & A worked example (stock 100, strike 100, 1 year, 5\% rate, 20\% volatility) has one correct call price, put price and vega & miss $<10^{-5}$\\
2 & Symbolic identities & The pricing formula must satisfy the Black-Scholes equation and put-call parity; MATLAB tests both & residual $<10^{-8}$\\
3 & Round trip & Price $\to$ volatility must return the volatility put in & error $<10^{-6}$\\
4 & Impossible prices & A call worth less than intrinsic value or more than the stock must be rejected, not ``solved'' & NaN returned\\
5 & Planted smile & Build a fake chain from a smile we choose, run the real pipeline, and confirm it recovers that smile at every quoted strike & error $<10^{-5}$\\
6 & Unseen points & Query the surface at strikes and maturities that were \emph{not} quoted and compare with the planted truth & error $<5\times10^{-3}$\\
7 & Real snapshot & On live data, MATLAB's volatilities must match the independent OCaml engine & max diff $<10^{-6}$\\\bottomrule
\end{tabularx}

Checks 5 and 6 are the strongest offline evidence. The planted smile is $\sigma(x)=0.25+0.6\,x^2-0.15\,x$ with $x=\ln(K/S)$, a curved, tilted shape like a real equity smile, quoted at three expiries (30, 60 and 120 days). Because the answer is known exactly, any error is the tool's own.

\section{Reference values, computed independently}
The expected numbers below were calculated outside MATLAB, with SciPy, so they do not share any code with the tool.

\begin{center}
\begin{tabular}{@{}lrr@{}}
\toprule
\textbf{Textbook case (S=K=100, T=1, r=5\%, $\sigma$=20\%)} & \textbf{Expected} & \textbf{Dividend yield 2\%}\\\midrule
Call price & 10.450584 & 9.227006\\
Put price & 5.573526 & 6.330081\\
Vega & 37.524035 & 37.901158\\
Put-call parity residual & 0 & 0\\\bottomrule
\end{tabular}
\end{center}

\section{Evidence gathered so far}
\begin{tabularx}{\linewidth}{@{}>{\raggedright\arraybackslash}p{0.34\linewidth}X>{\raggedright\arraybackslash}p{0.2\linewidth}@{}}
\toprule
\textbf{Evidence} & \textbf{Result} & \textbf{Where produced}\\\midrule
Independent Python solver (Brent root finder) vs OCaml on the AAPL snapshot of 12 Oct 2026 & 140 out-of-the-money contracts across 8 expiries, all solved; median difference $2.55\times10^{-7}$, maximum $4.99\times10^{-7}$ & Python, measured while writing this document\\
MATLAB pipeline vs OCaml on the same snapshot & Same 140 contracts; median $2.6\times10^{-7}$, maximum $5.0\times10^{-7}$; no bisection fallback needed & MATLAB, earlier session\\
Round-trip test grid (\texttt{test\_vol\_pipeline}) & 142 cases tested, 18 skipped where price carries no information about volatility; worst error $3.1\times10^{-7}$ & MATLAB, earlier session\\
Planted-smile surface, unseen points (check 6) & Largest miss about $1.1\times10^{-3}$ (0.1 volatility points), mean $1.5\times10^{-4}$ across 2000 random points & Python replica of the interpolation, not MATLAB\\\bottomrule
\end{tabularx}

\textbf{Reading the agreement.} The OCaml engine prints volatilities to six decimals, so the largest possible rounding difference is $5\times10^{-7}$. Both independent programs land exactly at that ceiling, which means the remaining difference is the printing, not the mathematics. Three unrelated implementations (OCaml, MATLAB, SciPy) agree on 140 real contracts.

\section{Procedure, step by step}
\begin{enumerate}[leftmargin=1.5em]
\item Open MATLAB in the repository's \texttt{matlab} folder.
\item Run \texttt{test\_vol\_pipeline} and \texttt{test\_surface}. Each prints \texttt{[ok]} lines and ends with ``All \ldots checks passed.''
\item Run \texttt{verify\_demo}. A report like the one below appears, followed by a figure comparing the planted smile with the recovered one at a 90-day maturity that was never quoted.
\item Run \texttt{verify\_demo("snapshots/<file>.csv")} on a saved snapshot to add check 7.
\item Run \texttt{verify\_independent.py} on the same snapshot and confirm the Python numbers sit near the MATLAB ones.
\item Record the outcome in the sign-off sheet below.
\end{enumerate}

\textbf{Shape of the report} (layout shown; numbers will be your own):
\begin{lstlisting}
=========== VERIFICATION REPORT ===========
[PASS] Textbook call, put and vega     call 10.450584, put 5.573526, ...
[PASS] Black-Scholes PDE holds         largest residual ...
[PASS] Put-call parity holds           largest residual ...
[PASS] Price to volatility round trip  ... worst error ...
[PASS] Impossible prices return NaN    both rejected
[PASS] Planted smile recovered at ...  ... worst error ...
[PASS] Surface accurate at points ...  400 points, worst error ...
-------------------------------------------
ALL 7 CHECKS PASSED
\end{lstlisting}

\section{If a check fails}
\begin{tabularx}{\linewidth}{@{}>{\raggedright\arraybackslash}p{0.27\linewidth}X@{}}
\toprule
\textbf{Failing check} & \textbf{Most likely cause}\\\midrule
1 or 2 & The pricing formula or its symbolic derivation changed (a sign or discount-factor error)\\
3 & The solver's bracket or tolerance changed, or the vega used for the Newton step is wrong\\
4 & The no-arbitrage guard was removed, so impossible prices produce numbers\\
5 & The forward-price estimate or the out-of-the-money selection is wrong, shifting strikes\\
6 & The interpolation changed: scaling of axes, or interpolating the wrong quantity\\
7 & Bad data in the snapshot, or the dividend-yield or interest-rate column differs between engines\\\bottomrule
\end{tabularx}

\section{Sign-off sheet}
\begin{tabularx}{\linewidth}{@{}>{\raggedright\arraybackslash}p{0.2\linewidth}XXXX@{}}
\toprule
\textbf{Date / ticker} & \textbf{test\_vol\_pipeline} & \textbf{test\_surface} & \textbf{verify\_demo} & \textbf{Python check}\\\midrule
\rule{0pt}{2.2em} & & & & \\
\rule{0pt}{2.2em} & & & & \\
\rule{0pt}{2.2em} & & & & \\\bottomrule
\end{tabularx}

\section{What this does not prove}
\begin{itemize}[leftmargin=1.5em]
\item That Yahoo's quotes are correct or current. The procedure verifies the \emph{computation}, not the data.
\item That the surface is free of arbitrage. The interpolated surface is a faithful picture of the quotes; a no-arbitrage fit is a planned later project.
\item That Black-Scholes is the right model. Its disagreement with the market is precisely what the surface displays.
\end{itemize}
\end{document}
